% généré par scripts/06_textes.py
\dosmeta{Année de classement : 382 (directe : nom consulaire lu). Références : ICUR I 1439 = ILCV 4700. Rattachement selon l’EDB : \emph{Incertae originis} (le volume I des ICVR réunit les inscriptions dont la provenance est perdue). Conservation selon l’EDB : \emph{Musei Vaticani: Lapidario Cristiano}. Support : Plaque de marbre. 30 × 58 cm (h. × l.). Inscription gravée. Fiche EDB : Filippo Piazzolla, première version 2015. Notice du catalogue p.~\pageref{not:231}.}
\txtbloc{Texte transmis (EDCS, état 2022)}{\lnum{1}Innocens recessit innocentium misericors fuit\par
\lnum{2}quis non doluit etati tuae piasq(ue) lacrimas fudit\par
\lnum{3}in te spes futura ex(s)pectabatur \{i per te\}\par
\lnum{4}per te ⸢g⸣loria perennis Celerine fili\par
\lnum{5}fidelis quiescis in pace ⸢q⸣ui ⸢v⸣ixi(t) ann(um) I\par
\lnum{6}m(enses) IIII d(epositus) VI Kal(endas) Oct(obres) Ant⸢o⸣n(i)o e⸢t⸣ Syagrio}
\txtbloc{Texte de l’EDB (EDB31204)}{\lnum{1}innocens recessit innocen̂tium mi`s´eric`o´rs fui`t´\par
\lnum{2}quis non(as) doluit ⟨a⟩etati tuae piasq(ue) lacrima`s´ fud`it´\par
\lnum{3}in te spes futura expectabatur 〚i〛 per te\par
\lnum{4}per te ⌜g⌝loria perennis Celerine fili\par
\lnum{5}fidelis quiescis in pace ⌜q⌝ui ⟨v⟩i\{i\}xit ann(is) I\par
\lnum{6}m(ensibus) IIII d(iebus) VI kal(endas) oct(obres) Antino e⌜t⌝ Syagrio}
